\section{Обзор существующих решений}
\label{sec:review}

В настоящем разделе рассмотрены существующие подходы к решению
поставленной задачи. Обзор структурирован по трём независимым
составляющим: методы локального планирования движения, методы
построения карты и локализации (SLAM) и программные навигационные
стеки, на базе которых практически реализуются такие решения. В конце
раздела выполнено обоснование выбранного подхода.

\subsection{Подходы к локальному планированию траектории}

\subsubsection{Реактивные методы}

Исторически первыми практическими подходами были реактивные локальные
планировщики: Dynamic Window Approach (DWA), Timed Elastic Band (TEB),
методы на основе потенциальных полей. DWA выбирает команду скорости из
допустимого по динамике окна, максимизирующую заданную эвристику. Такие подходы обладают
высокой вычислительной эффективностью, однако с ростом числа разнородных
требований эвристика становится громоздкой и плохо поддаётся настройке,
что ограничивает их применимость в задаче сопровождения.

\subsubsection{Классическое модельное прогнозирующее управление (MPC)}

Модельное прогнозирующее управление~\cite{Alkhaddad2023MPC} представляет
собой класс методов оптимального управления, в котором на каждом шаге
решается задача оптимизации с конечным горизонтом и использованием
предсказательной модели объекта и составной функции потерь. MPC
естественным образом формализует произвольное число разнородных
требований в виде слагаемых функционала и ограничений, а также легко
учитывает физические ограничения на управление.

Применение классического нелинейного MPC (например, с решателями на
основе SQP~\cite{Alkhaddad2023MPC}) требует аналитической гладкости
модели и штрафов, а также значительных вычислительных ресурсов для
поиска глобального минимума. Для роботов с колёсами Mecanum подходы
MPC рассмотрены в работе~\cite{MorenoCaireta2020MPC}, где учитываются
как кинематика, так и электрическая динамика приводов, но вычислительная
сложность остаётся высокой.

\subsubsection{Стохастические методы: MPPI}

Model Predictive Path Integral (MPPI) --- выборка-ориентированная
стохастическая реализация MPC. В каждом такте MPPI формирует набор
сэмплов управления путём добавления гауссовского шума к номинальному
управлению, прогнозирует траектории с использованием модели системы,
вычисляет стоимость каждой траектории и формирует обновление
номинального управления как экспоненциально взвешенное среднее по
сэмплам.

Преимущества MPPI перед классическим MPC в контексте задачи
сопровождения пользователя:
\begin{itemize}
    \item отсутствие требования гладкости функции потерь и модели
          системы, что позволяет использовать штрафы по данным costmap и
          произвольные логические условия (например, индикатор видимости
          цели);
    \item естественная параллелизация по независимым траекториям,
          позволяющая эффективно использовать современные CPU/GPU;
    \item устойчивость к локальным минимумам за счёт случайного
          сэмплирования.
\end{itemize}

Недостатки MPPI: необходимость тонкой настройки параметров
сэмплирования (число сэмплов, дисперсия шума, температура) и
чувствительность качества решения к количеству сэмплов.

\subsection{Методы SLAM для замкнутых помещений}

В закрытых помещениях, где отсутствует надёжный сигнал глобального
позиционирования, навигация строится на методах SLAM. Для 2D-лидара
наиболее распространены:
\begin{itemize}
    \item \textbf{Gmapping}~--- частично-фильтровое решение Rao-Blackwellized
          particle filter SLAM;
    \item \textbf{Hector~SLAM}~--- сопоставление сканов на основе
          градиентной оптимизации, не требующее одометрии;
    \item \textbf{Cartographer}~--- решение Google с графовой оптимизацией
          и замыканием цикла;
    \item \textbf{slam\_toolbox}~--- прямой аналог Cartographer в
          экосистеме ROS~2 с поддержкой асинхронного online-режима.
\end{itemize}

Сравнительные исследования~\cite{Sreerambatla2021SLAM,Nguyen2022SLAM}
показывают, что для задач квартирной навигации с 2D-лидаром решения
класса slam\_toolbox и Cartographer обеспечивают лучший баланс
точности карты, стабильности локализации и вычислительной нагрузки.
Hector~SLAM выигрывает в простоте запуска и отсутствии требования
одометрии, но проигрывает в точности на длинных траекториях из-за
отсутствия замыкания циклов.

\subsection{Навигационные стеки}

Практически значимые решения строятся не с нуля, а на основе
готовых навигационных стеков, обеспечивающих интеграцию планировщиков,
контроллеров, поведенческих деревьев и средств преобразования систем координат.

В среде ROS~2 де-факто стандартом является стек \textbf{Navigation~2}
(Nav2): он включает сервер глобального планировщика, сервер локального
контроллера с набором плагинов, сервер поведений (spin, wait, backup) и
BT~Navigator, выполняющий поведенческое дерево. Среди локальных
планировщиков Nav2 имеются реализации DWB, TEB, Regulated Pure Pursuit и
MPPI (\texttt{nav2\_mppi\_controller}), что делает Nav2 удобной базой для
задачи настоящей работы.

\subsection{Сравнительный анализ и обоснование выбора}

Сопоставление рассмотренных подходов с требованиями
главы~\ref{sec:problem} приводит к следующим выводам.
\begin{enumerate}
    \item Реактивные методы (DWA, потенциальные поля) не позволяют
          системно формализовать одновременно требования по удержанию
          дистанции и центрированию цели, поэтому не выбираются.
    \item Классический нелинейный MPC пригоден для формулировки задачи,
          но требует гладких штрафов, что затрудняет использование
          готовых инструментов Nav2 и costmap.
    \item \textbf{MPPI} поддерживает произвольные негладкие штрафы,
          естественным образом реализуется через архитектуру критиков,
          легко интегрируется в Nav2, имеет открытую реализацию для
          омни-платформ.
    \item Для построения карты и локализации выбрано решение
          \textbf{slam\_toolbox} как современный, поддерживаемый аналог
          Cartographer в экосистеме ROS~2.
    \item Для интеграции компонент выбран стек \textbf{Nav2}, что
          упрощает построение поведенческого слоя и использует готовую
          реализацию MPPI.
\end{enumerate}

Ни одно из существующих готовых решений не покрывает полностью
требования задачи в части одновременного удержания дистанции,
центрирования цели и обхода подвижных препятствий с устойчивым
поведением при потере цели, поэтому в рамках ВКР в состав MPPI-контроллера
добавлены пользовательские критики, формализующие соответствующие
требования (подробно описано в главе~\ref{sec:implementation}).
